% vim: set spell spelllang=en tw=100 et sw=4 sts=4 foldmethod=marker foldmarker={{{,}}} :

\documentclass{beamer}

\usepackage{tikz}
\usepackage{xcolor}
\usepackage{complexity}
\usepackage{hyperref}
\usepackage{microtype}
\usepackage{amsmath}                   % \operatorname
\usepackage{amsfonts}                  % \mathcal
\usepackage{amssymb}                   % \nexists
\usepackage{gnuplot-lua-tikz}          % graphs
\usepackage[vlined]{algorithm2e} % algorithms
\usepackage{centernot}
\usepackage{mathtools}
\usepackage{listings}

\usetikzlibrary{shapes, arrows, shadows, calc, positioning, fit}
\usetikzlibrary{decorations.pathreplacing, decorations.pathmorphing, shapes.misc}
\usetikzlibrary{tikzmark}

\definecolor{uofguniversityblue}{rgb}{0, 0.219608, 0.396078}

\definecolor{uofgheather}{rgb}{0.356863, 0.32549, 0.490196}
\definecolor{uofgaquamarine}{rgb}{0.603922, 0.72549, 0.678431}
\definecolor{uofgslate}{rgb}{0.309804, 0.34902, 0.380392}
\definecolor{uofgrose}{rgb}{0.823529, 0.470588, 0.709804}
\definecolor{uofgmocha}{rgb}{0.709804, 0.564706, 0.47451}
\definecolor{uofgsandstone}{rgb}{0.321569, 0.278431, 0.231373}
\definecolor{uofgforest}{rgb}{0, 0.2, 0.129412}
\definecolor{uofglawn}{rgb}{0.517647, 0.741176, 0}
\definecolor{uofgcobalt}{rgb}{0, 0.615686, 0.92549}
\definecolor{uofgturquoise}{rgb}{0, 0.709804, 0.819608}
\definecolor{uofgsunshine}{rgb}{1.0, 0.862745, 0.211765}
\definecolor{uofgpumpkin}{rgb}{1.0, 0.72549, 0.282353}
\definecolor{uofgthistle}{rgb}{0.584314, 0.070588, 0.447059}
\definecolor{uofgrust}{rgb}{0.603922, 0.227451, 0.023529}
\definecolor{uofgburgundy}{rgb}{0.490196, 0.133333, 0.223529}
\definecolor{uofgpillarbox}{rgb}{0.701961, 0.047059, 0}
\definecolor{uofglavendar}{rgb}{0.356863, 0.301961, 0.580392}

\tikzset{vertex/.style={draw, circle, inner sep=0pt, minimum size=0.5cm, font=\small\bfseries}}
\tikzset{notvertex/.style={vertex, color=white, text=black}}
\tikzset{plainvertex/.style={vertex}}
\tikzset{vertexc1/.style={vertex, fill=uofgburgundy, text=white}}
\tikzset{vertexc2/.style={vertex, fill=uofgsandstone, text=white}}
\tikzset{vertexc3/.style={vertex, fill=uofgforest, text=white}}
\tikzset{vertexc4/.style={vertex, fill=uofgheather, text=white}}
\tikzset{edge/.style={color=black!50!white}}
\tikzset{bedge/.style={ultra thick}}
\tikzset{edged/.style={color=screengrey, dashed}}
\tikzset{edgel3/.style={color=uofgrose, ultra thick}}

% {{{ theme things
\useoutertheme[footline=authortitle]{miniframes}
\useinnertheme{rectangles}

\setbeamerfont{block title}{size={}}
\setbeamerfont{title}{size=\large,series=\bfseries}
\setbeamerfont{section title}{size=\large,series=\mdseries}
\setbeamerfont{author}{size=\normalsize,series=\mdseries}
\setbeamercolor*{structure}{fg=uofguniversityblue}
\setbeamercolor*{palette primary}{use=structure,fg=black,bg=white}
\setbeamercolor*{palette secondary}{use=structure,fg=white,bg=uofgcobalt}
\setbeamercolor*{palette tertiary}{use=structure,fg=white,bg=uofguniversityblue}
\setbeamercolor*{palette quaternary}{fg=white,bg=black}

\setbeamercolor*{titlelike}{parent=palette primary}

\beamertemplatenavigationsymbolsempty

\setbeamertemplate{title page}
{
    \begin{tikzpicture}[remember picture, overlay]
        \node at (current page.north west) {
            \begin{tikzpicture}[remember picture, overlay]
                \fill [fill=uofguniversityblue, anchor=north west] (0, 0) rectangle (\paperwidth, -2.6cm);
            \end{tikzpicture}
        };

        \node (logo) [anchor=north east, shift={(-0.6cm,-0.6cm)}] at (current page.north east) {
            \includegraphics[keepaspectratio=true,scale=0.7]{UoG_keyline.pdf}
        };

        \node [anchor=west, xshift=0.2cm] at (current page.west |- logo.west) {
            \begin{minipage}{0.65\paperwidth}\raggedright
                {\usebeamerfont{title}\usebeamercolor[white]{}\inserttitle}\\[0.1cm]
                {\usebeamerfont{author}\usebeamercolor[white]{}\insertauthor}
            \end{minipage}
        };
    \end{tikzpicture}
}

\setbeamertemplate{section page}
{
    \begin{centering}
        \begin{beamercolorbox}[sep=12pt,center]{part title}
            \usebeamerfont{section title}\insertsection\par
        \end{beamercolorbox}
    \end{centering}
}

\newcommand{\frameofframes}{/}
\newcommand{\setframeofframes}[1]{\renewcommand{\frameofframes}{#1}}

\makeatletter
\setbeamertemplate{footline}
{%
    \begin{beamercolorbox}[colsep=1.5pt]{upper separation line foot}
    \end{beamercolorbox}
    \begin{beamercolorbox}[ht=2.5ex,dp=1.125ex,%
        leftskip=.3cm,rightskip=.3cm plus1fil]{author in head/foot}%
        \leavevmode{\usebeamerfont{author in head/foot}\insertshortauthor}%
        \hfill%
        {\usebeamerfont{institute in head/foot}\usebeamercolor[fg]{institute in head/foot}\insertshortinstitute}%
    \end{beamercolorbox}%
    \begin{beamercolorbox}[ht=2.5ex,dp=1.125ex,%
        leftskip=.3cm,rightskip=.3cm plus1fil]{title in head/foot}%
        {\usebeamerfont{title in head/foot}\insertshorttitle}%
        \hfill%
        {\usebeamerfont{frame number}\usebeamercolor[fg]{frame number}\insertframenumber~\frameofframes~\inserttotalframenumber}
    \end{beamercolorbox}%
    \begin{beamercolorbox}[colsep=1.5pt]{lower separation line foot}
    \end{beamercolorbox}
}

% }}}

\title{Linear Constraints}
\author{Ciaran McCreesh}

\begin{document}

{
    \usebackgroundtemplate{
        \tikz[overlay, remember picture]
        \node[at=(current page.south), anchor=south, inner sep=0pt]{\includegraphics[keepaspectratio=true, width=\paperwidth]{background4.jpg}};
    }
    \begin{frame}[plain,noframenumbering]
        \titlepage
    \end{frame}
}

\begin{frame}{Linear (or Scalar) Products}
    \begin{itemize}
        \item Let $c$ be an array of constants and $v$ an array of variables. The \emph{scalar
            product} of $c$ and $v$ is \begin{align*}
        c \cdot v = \sum_{i \in A}{c[i] \times v[i]}\end{align*}
        \item We constrain this with an operator to another variable (or constant). The operator can
            be an equation ($=$), or an inequality ($\le$ or $\ge$).
        \item In MiniZinc, the compiler automatically detects constraints like $3 * x + 2 * y <= 10$
            and $sum(i \operatorname{in} s)(a[i] * b[i]) = 12$.
        \item These are known as \emph{linear} equalities and inequalities, because
            geometrically they define a straight line.
    \end{itemize}
\end{frame}

\begin{frame}{What Can We Do?}
    \vspace*{-1em}
    \begin{align*}
        & x_1 \in \{ 0, 2 \} \\
        & x_2 \in \{ 0, 1 \} \\
        & x_3 \in \{ 0, 2 \} \\
        & 5x_1 + 2x_2 + x_3 \ge 8
    \end{align*}

    \begin{itemize}
        \item <2-> $x_1$ must be at least $1$, because we can get at most $(2 *
            1) + (1 * 2) = 4$ from the other two variables.
        \item <3-> $x_2$ and $x_3$ don't change, because $x_1 = 2$ can satisfy the inequality on its own.
        \item <4-> What if $x_1 \in \{ 0, 1 \}$?
            \begin{itemize}
            \item <5-> $x_2$ can't be $0$, because we can only get $(5 * 1) + (1 * 2) = 7$ from the other two
                without its help.
            \item <6-> $x_3$ can't be $0$, because we can only get $(5 * 1) + (2 * 1) = 7$ from the other two
                variables without its help.
            \end{itemize}
    \end{itemize}
\end{frame}

\begin{frame}{The Slack Algorithm for $\ge$ Inequalities}
    \begin{itemize}
        \item Feasibility: if every variable contributes as much as possible,
            do we have enough?
        \item Filtering: for each variable in turn, what's the most the
            remaining variables can contribute, and how much do I have to
            contribute if this happens?
            \begin{itemize}
                \item Assuming only positive numbers, this might increase lower bounds.
            \end{itemize}
        \item Actual algorithm removed for mental health reasons\ldots
    \end{itemize}
\end{frame}

\begin{frame}[fragile]{In Choco}
    \only<1> {
        \lstinputlisting[language=Java, basicstyle=\scriptsize\ttfamily, keywordstyle=\color{uofgcobalt}]{code/Scalar-snippet.java}
    }
    \only<2> {
        \begin{minipage}{0.47\paperwidth}\lstinputlisting[basicstyle=\scriptsize\ttfamily]{code/Scalar-output.txt}
        \end{minipage}\begin{minipage}{0.44\paperwidth}\lstinputlisting[basicstyle=\scriptsize\ttfamily]{code/Scalar-output2.txt}
        \end{minipage}
    }
\end{frame}

\begin{frame}{Remember Generalised Arc Consistency?}
    \begin{itemize}
        \item A \emph{binary} constraint is arc consistent if we can pick any value from either of its
            two variables, and find a supporting value in the other variable.
        \item A global constraint is generalised arc consistent (GAC) if each value
            is present in at least one solution to the constraint.
        \item Enforcing GAC will only touch upper and lower bounds for
            linear inequalities.
            \begin{itemize}
                \item We could also define \emph{bounds} consistency.
            \end{itemize}
    \end{itemize}
\end{frame}

\begin{frame}{Linear Equalities}
    \begin{itemize}
        \item What if we have an equation, rather than an inequality?
            \begin{align*}x = \sum_{i \in A}{c[i] \times v[i]}\end{align*}
        \item <2-> The subset sum problem is \NP-complete\ldots
        \item <3-> We can enforce GAC on a $\le$ constraint and a $\ge$
            constraint simultaneously. This is \textbf{not} the same as enforcing GAC
            on the equality.
    \end{itemize}
\end{frame}

\begin{frame}[fragile]{In Choco}
    \only<1> {
        \lstinputlisting[language=Java, basicstyle=\scriptsize\ttfamily, keywordstyle=\color{uofgcobalt}]{code/SmallScalarEquality-snippet.java}
    }
    \only<2> {
        \lstinputlisting[basicstyle=\scriptsize\ttfamily]{code/SmallScalarEquality-output.txt}
    }
\end{frame}

\begin{frame}{In Choco (Attempt Two)}
    \only<1> {
        \lstinputlisting[language=Java, basicstyle=\scriptsize\ttfamily, keywordstyle=\color{uofgcobalt}]{code/BigScalarEquality-snippet.java}
    }
    \only<2> {
        \begin{minipage}{0.47\paperwidth}\lstinputlisting[basicstyle=\scriptsize\ttfamily]{code/BigScalarEquality-output.txt}
        \end{minipage}\begin{minipage}{0.44\paperwidth}\lstinputlisting[basicstyle=\scriptsize\ttfamily]{code/BigScalarEquality-output2.txt}
        \end{minipage}
    }
\end{frame}

\begin{frame}{I'm Telling You This Twice Because It's Important}
    \begin{itemize}
        \item Propagation and consistency both depend upon how exactly you represent constraints.
        \item Modelling languages and solvers will \emph{sometimes} help you out.
    \end{itemize}
\end{frame}

\begin{frame}{Linear Programming}
    \begin{itemize}
        \item If all your variables are real numbers\ldots
        \item And all of your constraints are linear constraints\ldots
        \item And if your objective function is a linear sum\ldots
        \item \ldots Then you should use a \emph{linear programming} solver instead. These run in
            low order polynomial time\footnote{Terms and conditions apply.}.
    \end{itemize}
\end{frame}

\begin{frame}{Mixed Integer Programming}
    \begin{itemize}
        \item Like linear programming, except some or all variables can be required to be integers.
        \item This makes it hard again.
        \item Mixed integer programming solvers work very differently to constraint programming
            solvers, and can be much faster for some problems.
        \item Turning constraints into linear inequalities can be difficult. MiniZinc can compile
            \emph{some} constraints to MIP automatically.
    \end{itemize}
\end{frame}

\begin{frame}{SAT Solving (Going Even Further Off Topic)}
    \begin{itemize}
        \item All variables are booleans, and all constraints are in CNF.
        \item Much more restricted input than CP, but much faster supporting algorithms.
        \item Solvers are driven by \emph{unit propagation} and \emph{clause learning}.
            \begin{itemize}
                \item Rather than backtracking on failure, SAT solvers analyse the cause of the
                    failure, and produce a new constraint forbidding the cause.
                \item Propagating CNF much faster than propagating constraints
                    (terms and conditions apply).
            \end{itemize}
        \item MiniZinc can compile to SAT.
        \item Impossible to get GAC on all-different from a polynomial-sized
            CNF encoding, though\ldots
    \end{itemize}
\end{frame}

\begin{frame}{Jess Probably Isn't Mean Enough To Ask This On A Test}

    Consider the following constraint:
    \begin{align*}
        2x_1 - 2x_2 + 4x_3 \ge 11
    \end{align*}
    Given $x_1 \in \{ 2, 3, 4, 5 \}$, $x_2 \in \{ 1, 2, 3, 50 \}$ and $x_4 \in \{ 0, 1, 2
    \}$, prove $x_2 \ne 50$. What else can we infer?

    Now suppose we add a second constraint,
    \begin{align*}
        2x_1 - 2x_2 + 4x_3 \le 11
    \end{align*}
    Why is enforcing bounds consistency on these two constraints \emph{not} the same as enforcing
    bounds consistency on this single equality constraint?
    \begin{align*}
        2x_1 - 2x_2 + 4x_3 = 11
    \end{align*}

\end{frame}

\begin{frame}[plain,noframenumbering]
    \begin{tikzpicture}[remember picture, overlay]
        \node at (current page.north west) {
            \begin{tikzpicture}[remember picture, overlay]
                \fill [fill=uofguniversityblue, anchor=north west] (0, 0) rectangle (\paperwidth, -1.7cm);
            \end{tikzpicture}
        };

        \node (logo) [anchor=north east, shift={(-0.3cm,-0.2cm)}] at (current page.north east) {
            \includegraphics[keepaspectratio=true,scale=0.55]{UoG_keyline.pdf}
        };
    \end{tikzpicture}
\end{frame}

\end{document}
